\documentclass[12pt]{article}
\usepackage[utf8]{inputenc}
\usepackage[T1]{fontenc}
\usepackage{hyperref}
\title{Tsakonian Language Resources for K-12 Curriculum Development}
\author{Compiled by Claude Sonnet 5 (claude-sonnet-5) for Andr\'as Kornai}
\date{}
\begin{document}
\maketitle

\section{Introduction}
Tsakonian (ISO 639-3 tsd) is a critically endangered Indo-European language of the Peloponnese in Greece. It is commonly described as a descendant of Doric Greek. Only older adults speak it as a first language, and it has no official status. Searches found efforts to teach it in local schools, but these do not seem to have had much success. Little curricular material exists. This report lists only what the searches turned up.

\section{Primary Educational Resources}
No dedicated textbook, primer or structured school course was found. One news report says some teaching materials were produced for local schools, but it names none and gives no URL. The digital lexicon by Panagiotis Tsagouris, headmaster of the Leonidio primary school, is the closest thing to a school resource. His students record the words, pronunciation and usage. It had about 4,650 Greek-translated entries, many with audio. See \url{https://proof.csmonitor.com/World/Making-a-difference/2016/0818/How-a-headmaster-is-trying-to-save-an-ancient-language} (also mirrored at \url{https://myhero.com/how-a-headmaster-is-trying-to-save-an-ancient-language}).

\section{Linguistic and Reference Materials}
\begin{itemize}
\item C.\ A.\ Scutt, ``The Tsakonian Dialect'' (Annual of the British School at Athens, 1912--1914). Glottolog catalogues it as a grammar sketch: \url{https://glottolog.org/resource/reference/id/318490}
\item General overview of Tsakonian Greek: \url{https://en.wikipedia.org/wiki/Tsakonian_Greek}
\item Endangered Languages Project entry for Tsakonian: \url{https://endangeredlanguages.com/node/106314/printable/print}. A search summary mentioned a Project Hom\`ere grammar resource listed there. This was not confirmed.
\end{itemize}
A 2024 article from St Petersburg State University describes an online Tsakonian--Greek dictionary with Greek-to-Tsakonian search. The authors suggest it could become a corpus and a teaching platform. The listing is at \url{https://pureportal.spbu.ru/en/publications/---------(a5e09c36-8aa8-4559-9b7f-79eeb475073c).html}. I did not verify that the dictionary itself is online.

\section{Religious and Community Resources}
Nothing was found in the searches.

\section{Digital and Audio Resources}
The Leonidio digital lexicon (see above) includes audio clips for many entries. No other audio or video collections were found.

\section{Curriculum Development Context}
There is no known school curriculum. The lexicon and the SPbU dictionary could serve as vocabulary bases. Any curriculum would have to be built largely from scratch and from elderly speakers. Local contacts in Leonidio would need to be approached directly.
\end{document}